\section{Special Rules}

Except where changed below, use all standard rules from sections 1.0--15.0.

\begin{enumerate}
  \item \textbf{Royalist initiative.} Skip the Allied Rally, Artillery, March,
    and Combat Phases on Turn 1. Begin play with the Royalist Rally Phase and
    complete the rest of the turn normally. Starting with Turn 2, use the normal
    sequence of play.

  \item \textbf{Daylight.} Visibility is Clear throughout the scenario. Skip
    every Visibility Phase. The flash-rain rule and the Short Game's temporary
    cancellation of ditch effects are not used; ditch hexsides function
    normally from the beginning.

  \item \textbf{Uncommitted arrivals.} A Foot unit still off-map at the end of the
    scenario awards the Royalist player VPs equal to half its ordered combat
    strength, rounded down. Each Allied artillery piece still off-map awards 5
    VPs. These pieces do not count toward Allied demoralization.

  \item \textbf{Optional rules.} For the first playtest, do not use the Optional
    Skirmishing Rule. It can be added once the basic arrival schedule has been
    tested; it gives the exposed Allied rearguard another way to slow the
    Royalist advance and may require moving one Scottish group a turn later.
\end{enumerate}

\section{Victory}

Calculate VPs normally under 15.0, including the standard awards for the train,
artillery, leaders, demoralization, eliminated units, and disrupted Foot. Add
any VPs for Allied pieces that never entered under the special rule above.

Apply the standard victory ratios, with one scenario-specific change:

\begin{redBox}{The Royalists Must Exploit the Opportunity}
If neither player has a 1.5-to-1 ratio at the end of Turn 12, the Allied player
wins a \textbf{marginal scenario victory}. The Allied army has survived the
dangerous hours, completed its deployment, and denied Rupert the quick decision
he sought. A Royalist marginal, major, or crushing victory therefore requires
the corresponding standard ratio.
\end{redBox}

If the Allied player earns a standard victory ratio, determine its level
normally. The scenario victory described above is always marginal.
